\section{Absurdity in Discussion}

\begin{quotex}
The most difficult thing of all is to believe in truth because it is true; i.e., because it demands humble recognition and self-abnegation; it is far easier to accept truth as my opinion which I posit as truth.
\flright{\textsc{Sergius Bulgakov}}

Man comes from God; sin from man; ignorance and error, like pain and death, from sin; fallibility from ignorance; and from fallibility, absurdity in discussion.
\flright{\textsc{Donoso Cortes}}
\end{quotex}


In contrast to the authentic man who is open to the unconcealment of Being or Truth, \textbf{Martin Heidegger} in his \emph{Being and Time} described the inauthentic man as one immersed in the everyday, and whose consciousness is subject to the ``They'', that is, the undifferentiated mass. The primary means of this is idle chit chat (\emph{Gerede}).

Instead of discourse as the means to communicate truth, gain understanding, or share factual knowledge, in idle chatter, it becomes, instead, the means to form alliances, cement friendships, and create a sort of tribal consensus. Hoary arguments, long ago discredited, are repeated as new revelations and historical inaccuracies or misconceptions are believed with passionate conviction. The Internet does much to enable this on a wider scale, hence, we witness new web sites cropping boasting of their presentation of divergent points of view and the ``lively'' discussions which then ensue. Rather than leading to an openness to Truth, such discussion enclose one in a web of words and opinions. One takes comfort in the verbal support of fellow tribesmen and fantasizes about the \emph{bon mot} that will intellectually crush his opponent.

Anyone with a longtime interest in Chess knows that the opening moves are all predictable. Most are fruitless and will lead to quick defeat. The good openings, and all their variations, are typically memorized by a chess master, even up to 15 to 20 moves. Only at that point can the game get interesting. The same applies to Internet discussions. Most arrive stillborn at the first sentence. The arguments that have some potential only make it to the first or second stage of argument and counter-argument, before reaching the point of incommunicability. Should those who know many more stages in the syllogism attempt to participate, they are drowned out and misunderstood, and, like Sisyphus, are forever condemned to start again at the beginning.

Ultimately, the problem is that the rational mind is incapable of reaching first principles, since its function is evaluation and judging ---for that, it requires raw material to process and mediate. First principles can only come from something higher --- the intellect, or \emph{nous}, wherein lies the opening for the unconcealment of truth so that it can be directly perceived.


\flrightit{Posted on 2010-04-03 by Cologero}

\begin{center}* * *\end{center}

\begin{footnotesize}
\begin{sffamily}
\texttt{GFJF on 2010-04-03 at 23:36 said:}

I'm starting to see what you mean about not being partisan. There's traditionalism, and then there's Tradition. The discussion recently sparked by Stephen McNallen's article is a case in point. Some people have left some curious comments which seem inexplicable unless they were fired by an anti-Christian animus.


\hfill

\end{sffamily}
\end{footnotesize}

